\documentclass[11pt,letterpaper]{article}
\PassOptionsToPackage{hyphens}{url}
\usepackage[margin=.83in,top=.8in,bottom=.8in,headheight=15pt]{geometry}
\usepackage[T1]{fontenc}
\usepackage{lmodern,amsmath,amssymb,mathtools,bm,booktabs,array,graphicx,tikz}
\usetikzlibrary{arrows.meta,positioning,calc,decorations.pathreplacing,patterns}
\definecolor{ink}{HTML}{17344D}\definecolor{accent}{HTML}{087E8B}\definecolor{pale}{HTML}{EDF5F5}\definecolor{gold}{HTML}{AD6B19}
\providecommand{\dd}{\,\mathrm{d}}\providecommand{\R}{\mathbb{R}}
\providecommand{\E}{\mathbb{E}}\providecommand{\Var}{\operatorname{Var}}
\providecommand{\Cov}{\operatorname{Cov}}\providecommand{\diag}{\operatorname{diag}}
\providecommand{\norm}[1]{\left\lVert#1\right\rVert}
\providecommand{\abs}[1]{\left|#1\right|}
\usepackage{microtype,hyperref,fancyhdr,enumitem,amsthm,sectsty,needspace}
\allsectionsfont{\raggedright}
\hypersetup{colorlinks=true,linkcolor=ink,urlcolor=accent,pdfauthor={Zan Ahmad},pdftitle={Circulatory models and parameter identifiability / Review}}
\setlength{\parindent}{0pt}\setlength{\parskip}{6pt}\linespread{1.07}
\setlist{itemsep=3pt,topsep=4pt,leftmargin=*}
\emergencystretch=2em\clubpenalty=10000\widowpenalty=10000
\pagestyle{fancy}\fancyhf{}\fancyhead[L]{\small Computational Cardiology}\fancyhead[R]{\small Lecture 1}
\fancyfoot[L]{\footnotesize Zan Ahmad}\fancyfoot[R]{\footnotesize\thepage}\renewcommand{\headrulewidth}{.3pt}
\newtheorem{proposition}{Proposition}\newtheorem{theorem}{Theorem}\newtheorem{lemma}{Lemma}
\theoremstyle{definition}\newtheorem{definition}{Definition}\newtheorem{example}{Example}
\numberwithin{equation}{section}
\begin{document}
\section*{Circulatory models and parameter identifiability --- review sheet}
Use this after attempting the worked examples and practice. Explain each claim in words before calculating.

\section*{Key ideas}
\begin{itemize}
\item Question and observable determine useful resolution.
\item The six-region boundary represents a closed circulation, not six anatomical heart chambers.
\item Mean flow is a product under the stated repeated-cycle assumption.
\item Equal outputs can arise from different hidden parameters.
\item Verification and validation answer different questions.
\end{itemize}
\section*{Common errors and repairs}
\textbf{Inferring a unique internal mechanism from one successful fit.} Look for transformations that leave the observable unchanged and identify additional measurements.

\textbf{Using a pressure gauge shift as if it were a change of transmural loading.} Track the exterior pressure and the storage relation too.

\section*{Retrieve without your notes}
\begin{itemize}
\item What converts rate times stroke volume from mL/min to L/min?
\item Can a more detailed model remove missing measurement information by itself?
\end{itemize}
\section*{Readiness check}
\begin{itemize}
\item Trace the two-pump circulation and justify a six-region lumped boundary.
\item Map three scientific questions to assumptions, model resolution and measurable outputs.
\item Calculate mean flow and effective resistance with compatible units.
\item Demonstrate a non-identifiable parameter pair and select an additional observation.
\item Distinguish numerical verification, parameter estimation and experimental validation.
\end{itemize}
\end{document}
